\documentclass[../differential_topology.tex]{subfiles}
\begin{document}
\section{Aula 15 - 05 de Novembro, 2026}
\subsection{Motivações}
\begin{itemize}
	\item Propriedades dos Graus
\end{itemize}
\subsection{Propriedades dos Graus}
Começamos provando os dois lemas utilizados na aula passada:
\begin{lemma*}
	Suponha que \(M = \partial X\), onde X é uma variedade orientável com bordo. Se \(f:M\rightarrow N\) admite uma extensão \(F:X\rightarrow N\), então \(\mathrm{deg}(f) = 0\).
\end{lemma*}
\begin{proof*}
	Se \(y\in N\) é um valor regular de F, então \(F^{-1}(y)\) é uma subvariedade compacta de dimensão 1, pois
	\[
		\mathrm{dim}_{}(F^{-1}(y)) = \mathrm{dim}_{}(X) - \mathrm{dim}_{}(N) = (n+1) - n = 1,
	\]
	que satisfaz
	\[
		\partial F^{-1}(y) = F^{-1}(y)\cap \partial M.
	\]
	Pela classificação das variedades compactas de dimensão 1, \(F^{-1}(y)\) é uma união finita e disjunta de circunferências e de arcos fechados, e como essas componentes não possuem bordo, cada arco possui exatamente dois pontos de extremidades. Considere A uma dessas componentes homeomorfas a um arco com extremidades \(a, b\) em M e orientemos ela por meio das orientação de X e N.

	Para cada ponto \(x\in A\), escolhemos uma base positivamente orientada de \(T_{x}X\), digamos
	\[
		(v_1(x), v_2(x), \dotsc , v_{n+1}(x)),
	\]
	de maneira que
	\[
		T_{x}A = \langle v_1(x) \rangle
	\]
	com base positiva
	\[
		\left(\mathrm{d}F_x(v_2(x)), \dotsc , \mathrm{d}F_{x}(v_{n+1}(x))\right)
	\]
	Por isso, \(v_1(x)\) determina a orientação de \(A\), e escolhemos a orientação do arco de modo que, na extremidade \(a\), o vetor tangente aponte para fora de X, enquanto que, em \(b\), aponte para dentro; pela orientação induzida no bordo, as duas extremidades contribuem com sinais opostos para o grau de \(f\) devido à convenção da normal exterior (faz com que o sinal muda quanto o vetor tangente orientado passa de fora para dentro!), tal que
	\[
		\mathrm{sgn}(\mathrm{d}f_{a}) + \mathrm{sgn}(\mathrm{d}f_{b}) = 0.
	\]

	Consequentemente, cada componente de \(F^{-1}(y)\) que é um arco contribui em nada para a soma que define o grau, ao passo que, como as componentes homeomorfas às circunferências não possuem pontos no bordo, também não contribuem. Portanto,
	\[
		\mathrm{deg}(f) = 0. \text{ \qedsymbol}
	\]
\end{proof*}
\begin{lemma*}
	Se \(f\) e \(g\) são aplicações isotópicas, então para todo valor regular comum, vale
	\[
		\mathrm{deg}(f, y) = \mathrm{deg}(g, y).
	\]
\end{lemma*}
\begin{proof*}
	Como \(f\) e \(g\) são isotópicas, existe uma aplicação suave
	\[
		F:M\times [0, 1]\rightarrow N
	\]
	com
	\[
		F(x, 0) = f(x), \quad F(x, 1) = g(x);
	\]
	Denotando-na por \(F_{t}(x) = F(x, t)\), podemos escrever simplesmente, para facilitar,
	\[
		F_{0} = f, \quad F_{1} = g.
	\]

	Dito isso, escolha \(y\in N\) que seja ao mesmo tempo um valor regular de \(F,\; f\) e de \(g\), a qual existe devido ao \hyperlink{sard_theorem}{\textit{Teorema de Sard}}. Tome a variedade com bordo \(X = [0, 1]\times M\), com a orientação usual do produto e tal que o bordo seja formado por duas cópias de M com orientações opostas:
	\[
		\partial X = M_{1} \sqcup -M_{0},
	\]
	onde usamos a estratégia super útil\footnote{Axioma da Escolha!!} de
	\[
		M_{i} = M\times i, \quad i = 0, 1.
	\]

	Pelo teorema da pré-imagem, \(F^{-1}(y)\) é uma subvariedade compacta de dimensão 1 com
	\[
		\partial F^{-1}(y) = F^{-1}(y)\cap (M_{0}\cup M_1),
	\]
	e assim como no Lema 1, as componentes de \(F^{-1}(y)\) são arcos e circunferências, com cada um possuindo duas extremidades e cujas contribuições para o grau se cancelam.

	Utilizando o Lema anterior à aplicação \(F\), concluímos que
	\[
		\deg{(F|_{\partial X})}=0;
	\]
	no entanto, a orientação de \(M_1\) coincide com a original de \(M\), ao passo que \(M_{0}\) é oposta e
	\[
		\mathrm{deg}(F|_{\partial X}) = \deg{(g)} - \deg{(f)}.
	\]
	Portanto,
	\[
		\deg{(g)} - \deg{(f)} = 0 \Longleftrightarrow \deg{(g)} = \deg{(f)}. \text{ \qedsymbol}
	\]
\end{proof*}
\begin{prop*}[Grau da composição]
	Sejam M, N e P variedades diferenciáveis, orientadas e de mesma dimensão, com M e N compactas e N e P conexas. Sejam ainda as aplicações suaves
	\[
		f:M\rightarrow N \quad\&\quad g:N\rightarrow P.
	\]
	Então,
	\[
		\mathrm{deg}(g\circ f) = \mathrm{deg}(g)\cdot\mathrm{deg}(f).
	\]
\end{prop*}
\begin{proof*}
	Escolha um ponto \(a\in P\) que seja valor regular tanto de g quanto de \(g\circ f\); tal escolha é possível pelo \hyperlink{sard_theorem}{\textit{Teorema de Sard}}. Pela regra da cadeia,
	\[
		(g\circ f)'(x) = g'(f(x))\circ f'(x).
	\]
	Como a diferencial da composição é um isomorfismo para cada \(x\in (g\circ f)^{-1}(a)\), e como os espaços tangentes possuem a mesma dimensão, as diferenciais \(g'(f(x))\) e \(f'(x)\) também são isomorfismos. Assim, todo \(y\in g^{-1}(a)\) é valor regular de \(f\) (inclusive quando \(f^{-1}(y)=\emptyset\)).

	Pela definição do grau e pela multiplicatividade do sinal do determinante, temos
	\begin{align*}
		\mathrm{deg}(g\circ f)
		 & = \sum\limits_{x\in (g\circ f)^{-1}(a)}\mathrm{sgn}\bigl((g\circ f)'(x)\bigr)                                                                                    \\
		 & = \sum\limits_{x\in (g\circ f)^{-1}(a)}\mathrm{sgn}\bigl(g'(f(x))\bigr)\cdot\mathrm{sgn}\bigl(f'(x)\bigr)                                                        \\
		 & = \sum\limits_{y\in g^{-1}(a)}\sum\limits_{x\in f^{-1}(y)}\mathrm{sgn}\bigl(g'(y)\bigr)\cdot\mathrm{sgn}\bigl(f'(x)\bigr)                                        \\
		 & = \sum\limits_{y\in g^{-1}(a)}\mathrm{sgn}\bigl(g'(y)\bigr)\underbrace{\left(\sum\limits_{x\in f^{-1}(y)}\mathrm{sgn}\bigl(f'(x)\bigr)\right)}_{\mathrm{deg}(f)} \\
		 & = \mathrm{deg}(f)\sum\limits_{y\in g^{-1}(a)}\mathrm{sgn}\bigl(g'(y)\bigr)                                                                                       \\
		 & = \mathrm{deg}(f)\cdot\mathrm{deg}(g). \text{ \qedsymbol}
	\end{align*}
\end{proof*}

\begin{example}[Aplicação antipodal]
	Para \(n\geq 1\), considere a aplicação antipodal \(A:\mathbb{S}^{n}\rightarrow\mathbb{S}^{n}\), dada por
	\[
		A(x_1,\dotsc,x_{n+1})=(-x_1,\dotsc,-x_{n+1}).
	\]
	Ela pode ser escrita como uma composição de \(n+1\) reflexões,
	\[
		A=R_1\circ R_2\circ\dotsc\circ R_{n+1},
	\]
	onde \(R_i\) é a reflexão definida em aulas passadas. Cada reflexão inverte a orientação da esfera, de modo que \(\mathrm{deg}(R_i)=-1\) e pela proposição anterior,
	\[
		\mathrm{deg}(A)=\prod\limits_{i=1}^{n+1}\mathrm{deg}(R_i)=(-1)^{n+1}.
	\]
	Como \(\mathrm{deg}(\mathrm{id}_{\mathbb{S}^{n}})=1\), a invariância do grau por homotopia implica que, se n é par,
	\[
		A\not\cong\mathrm{id}_{\mathbb{S}^{n}}.
	\]
	Quando n é ímpar, a antipodal é homotópica à identidade: nesse caso, \(n+1\) é par e podemos obter uma homotopia por meio de rotações simultâneas nos planos coordenados \((x_1,x_2),\dotsc,(x_n,x_{n+1})\).
\end{example}

\subsection{O Teorema da Bola Cabeluda e Outra Prova do TFA}
\begin{def*}
	Seja \(M\subset\mathbb{R}^{k}\) uma variedade diferenciável. Um \textbf{campo vetorial tangente suave} sobre M é uma aplicação suave
	\[
		v:M\rightarrow\mathbb{R}^{k}
	\]
	tal que \(v(x)\in T_xM\) para todo \(x\in M\). Dizemos que v é \textbf{não nulo} se \(v(x)\neq 0\) para todo \(x\in M\). \(\square\)
\end{def*}

\begin{example}
	Quando \(M=\mathbb{S}^{n}\subset\mathbb{R}^{n+1}\), temos
	\[
		T_x\mathbb{S}^{n}=\{w\in\mathbb{R}^{n+1}:\langle x,w\rangle=0\}.
	\]
	Assim, v é tangente à esfera se, e somente se, \(\langle x,v(x)\rangle=0\) para todo \(x\in\mathbb{S}^{n}\).
\end{example}

\hypertarget{hairy_ball}{\begin{theorem*}[Teorema da Bola Cabeluda]
		A esfera \(\mathbb{S}^{n}\) admite um campo vetorial tangente suave e não nulo se, e somente se, n é ímpar.
	\end{theorem*}}
\begin{proof*}
	Suponha, inicialmente, que exista um campo tangente suave \(v:\mathbb{S}^{n}\rightarrow\mathbb{R}^{n+1}\), com \(v(x)\neq0\) para todo x. Dividindo v por sua norma, podemos supor \(\Vert v(x)\Vert=1\) em toda a esfera.

	Defina a aplicação
	\[
		F:\mathbb{S}^{n}\times[0,\pi]\rightarrow\mathbb{S}^{n},\quad F(x,\theta)=x\cos\theta+v(x)\sin\theta.
	\]
	Como v é tangente à esfera, vale \(\langle x,v(x)\rangle=0\); logo,
	\begin{align*}
		\Vert F(x,\theta)\Vert^{2}
		 & =\Vert x\Vert^{2}\cos^{2}\theta+2\langle x,v(x)\rangle\cos\theta\sin\theta+\Vert v(x)\Vert^{2}\sin^{2}\theta \\
		 & =\cos^{2}\theta+\sin^{2}\theta=1.
	\end{align*}
	Logo, F está bem definida e fornece uma homotopia entre a identidade e a antipodal, pois
	\[
		F(x,0)=x \quad\&\quad F(x,\pi)=-x.
	\]
	Pela invariância do grau por homotopia e pelo exemplo anterior,
	\[
		1=\mathrm{deg}(\mathrm{id}_{\mathbb{S}^{n}})=\mathrm{deg}(A)=(-1)^{n+1},
	\]
	o que implica que n é ímpar.

	Reciprocamente, se n é ímpar, \(n+1\) é par. Defina
	\[
		v(x_1,x_2,\dotsc,x_n,x_{n+1})=(-x_2,x_1,\dotsc,-x_{n+1},x_n).
	\]
	Temos \(\langle x,v(x)\rangle=0\) e \(\Vert v(x)\Vert=\Vert x\Vert=1\). Portanto, v é um campo tangente suave e não nulo sobre \(\mathbb{S}^{n}\), demonstrando a recíproca. \qedsymbol
\end{proof*}

\begin{tcolorbox}[
		skin=enhanced,
		title=Observação,
		fonttitle=\bfseries,
		colframe=black,
		colbacktitle=cyan!75!white,
		colback=cyan!15,
		colbacklower=black,
		coltitle=black,
		drop fuzzy shadow,
		%drop large lifted shadow
	]
	Se \(f:M\rightarrow N\) é uma aplicação suave entre variedades orientadas de mesma dimensão, com M compacta e N conexa, e f não é sobrejetora, então existe \(c\in N\setminus f(M)\). Como \(f^{-1}(c)=\emptyset\), c é um valor regular de f e
	\[
		\mathrm{deg}(f)=\sum\limits_{x\in f^{-1}(c)}\mathrm{sgn}(f'(x))=0.
	\]
	Em particular, vale a contrapositiva
	\[
		\mathrm{deg}(f)\neq0\Rightarrow f\text{ é sobrejetora}.
	\]
\end{tcolorbox}

\hypertarget{fundamental_theorem_algebra}{\begin{theorem*}[Teorema Fundamental da Álgebra]
		Todo polinômio não constante com coeficientes complexos possui pelo menos uma raiz complexa.
	\end{theorem*}}
\begin{proof*}
	Considere um polinômio mônico de grau \(n\geq1\), escrito como
	\[
		p(z)=z^{n}+c_{n-1}z^{n-1}+\dotsc+c_1z+c_0=z^{n}+Q(z),
	\]
	onde \(\mathrm{deg}(Q)<n\), quando \(Q\neq0\). Consideremos a homotopia
	\[
		F:\mathbb{C}\times[0,1]\rightarrow\mathbb{C},\quad F(z,t)=z^{n}+tQ(z),
	\]
	de modo que \(F_0(z)=z^{n}\) e \(F_1(z)=p(z)\).

	Como Q possui grau menor que n, podemos escolher \(R>0\) suficientemente grande para que
	\[
		|Q(z)|<|z|^{n},\quad\text{se }|z|=R.
	\]
	Então, para todo \(t\in[0,1]\) e \(|z|=R\), temos \(|tQ(z)|<|z|^{n}\), levando a \(F(z,t)\neq0\). Isso permite normalizar F sobre a circunferência e definir
	\[
		H:\mathbb{S}^{1}\times[0,1]\rightarrow\mathbb{S}^{1},\quad
		H(u,t)=\frac{F(Ru,t)}{|F(Ru,t)|};
	\]
	notemos que
	\[
		H(u,0)=u^{n} \quad\&\quad H(u,1)=\frac{p(Ru)}{|p(Ru)|}.
	\]
	Como H é uma homotopia e a aplicação \(u\mapsto u^{n}\) possui grau n (cada valor regular tem n pré-imagens com sinal positivo), segue que
	\[
		\mathrm{deg}(H_1)=\mathrm{deg}(H_0)=n\neq0.
	\]

	Suponha, por absurdo, que p não possua raízes complexas; então, a aplicação suave
	\[
		G:D_R\rightarrow\mathbb{S}^{1},\quad G(z)=\frac{p(z)}{|p(z)|},\quad
		D_R=\{z\in\mathbb{C}:|z|\leq R\},
	\]
	estaria bem definida em todo o disco fechado e satisfaria \(G(Ru)=H_1(u)\) em seu bordo. Logo, \(H_1\) se estenderia ao disco e seria homotópica a uma aplicação constante: por exemplo, por meio de
	\[
		K(u,t)=G((1-t)Ru),\quad t\in[0,1].
	\]
	Pela invariância do grau por homotopia, teríamos \(\mathrm{deg}(H_1)=0\), contradizendo \(\mathrm{deg}(H_1)=n\neq0\). Portanto, existe \(z_0\in\mathbb{C}\) tal que
	\[
		p(z_0)=0. \text{ \qedsymbol}
	\]
\end{proof*}

\end{document}
